\documentclass[12pt,a4paper]{article}
\usepackage[UTF8]{ctex}
\usepackage{amsmath,amssymb,amsthm}
\usepackage{geometry}
\usepackage{graphicx}
\usepackage{hyperref}
\usepackage{bm}

\geometry{margin=2.5cm}

\title{正交矩阵的性质详解}
\author{Modern Robotics}
\date{}

\begin{document}

\maketitle

\tableofcontents
\newpage

\section{引言}

正交矩阵是线性代数中非常重要的概念，在机器人学中有着广泛的应用，特别是在旋转矩阵、坐标变换等方面。本文档详细解释正交矩阵的定义、性质，特别是关于向量长度保持的性质。

\section{正交矩阵的定义}

\subsection{基本定义}

\textbf{定义1}：对于方阵$Q \in \mathbb{R}^{n \times n}$，如果满足
\begin{equation}
Q^T Q = I
\label{eq:orthogonal_def1}
\end{equation}

则称$Q$是\textbf{正交矩阵}。

\textbf{定义2}（等价定义）：如果$Q$满足
\begin{equation}
Q Q^T = I
\label{eq:orthogonal_def2}
\end{equation}

则$Q$也是正交矩阵。

\textbf{重要结论}：对于方阵，$Q^T Q = I$当且仅当$Q Q^T = I$。

\textbf{证明}：如果$Q^T Q = I$，则$Q$可逆，且$Q^{-1} = Q^T$。因此：
\begin{equation}
Q Q^T = Q Q^{-1} = I
\end{equation}

反之亦然。

\subsection{单位正交矩阵（标准正交矩阵）}

\textbf{定义}：如果正交矩阵$Q$的列向量（或行向量）不仅是正交的，而且是\textbf{单位向量}（长度为1），则称$Q$是\textbf{单位正交矩阵}或\textbf{标准正交矩阵}。

\textbf{等价表述}：
\begin{itemize}
\item $Q$的列向量是标准正交的：$\bm{q}_i^T \bm{q}_j = \delta_{ij}$（$\delta_{ij}$是Kronecker delta）
\item $Q$的行向量是标准正交的：$\bm{q}_i \bm{q}_j^T = \delta_{ij}$
\end{itemize}

\textbf{注意}：实际上，所有正交矩阵的列向量和行向量都是单位向量，所以"单位正交矩阵"和"正交矩阵"通常指的是同一个概念。但在某些文献中，"单位正交矩阵"强调列向量的单位长度性质。

\section{正交矩阵的基本性质}

\subsection{性质1：逆矩阵等于转置}

\textbf{性质}：如果$Q$是正交矩阵，则
\begin{equation}
Q^{-1} = Q^T
\label{eq:inverse_transpose}
\end{equation}

\textbf{证明}：由定义$Q^T Q = I$，直接得到$Q^{-1} = Q^T$。

\textbf{推论}：
\begin{itemize}
\item $(Q^T)^{-1} = Q$
\item $Q^T$也是正交矩阵
\end{itemize}

\subsection{性质2：行列式的值}

\textbf{性质}：如果$Q$是正交矩阵，则
\begin{equation}
\det(Q) = \pm 1
\label{eq:determinant}
\end{equation}

\textbf{证明}：
\begin{align}
\det(Q^T Q) &= \det(I) = 1 \\
\det(Q^T) \det(Q) &= 1 \\
(\det(Q))^2 &= 1 \\
\det(Q) &= \pm 1
\end{align}

\textbf{分类}：
\begin{itemize}
\item 如果$\det(Q) = 1$，则$Q$是\textbf{特殊正交矩阵}（旋转矩阵）
\item 如果$\det(Q) = -1$，则$Q$是\textbf{反射矩阵}（包含反射的旋转）
\end{itemize}

\subsection{性质3：保持向量长度（最重要的性质）}

\textbf{性质}：如果$Q$是正交矩阵，$\bm{y} \in \mathbb{R}^n$是任意向量，则
\begin{equation}
\|Q\bm{y}\| = \|\bm{y}\|
\label{eq:preserve_norm}
\end{equation}

\textbf{证明}：
\begin{align}
\|Q\bm{y}\|^2 &= (Q\bm{y})^T (Q\bm{y}) \\
&= \bm{y}^T Q^T Q \bm{y} \\
&= \bm{y}^T I \bm{y} \\
&= \bm{y}^T \bm{y} \\
&= \|\bm{y}\|^2
\end{align}

因此：$\|Q\bm{y}\| = \|\bm{y}\|$。

\textbf{几何意义}：正交矩阵保持向量的长度不变，因此表示\textbf{等距变换}（isometry）。

\subsection{性质4：保持内积}

\textbf{性质}：如果$Q$是正交矩阵，$\bm{x}, \bm{y} \in \mathbb{R}^n$是任意向量，则
\begin{equation}
(Q\bm{x})^T (Q\bm{y}) = \bm{x}^T \bm{y}
\label{eq:preserve_inner_product}
\end{equation}

\textbf{证明}：
\begin{align}
(Q\bm{x})^T (Q\bm{y}) &= \bm{x}^T Q^T Q \bm{y} \\
&= \bm{x}^T I \bm{y} \\
&= \bm{x}^T \bm{y}
\end{align}

\textbf{几何意义}：正交矩阵保持向量之间的夹角不变。

\textbf{推论}：由于$\cos\theta = \frac{\bm{x}^T \bm{y}}{\|\bm{x}\| \|\bm{y}\|}$，且$\|Q\bm{x}\| = \|\bm{x}\|$，$\|Q\bm{y}\| = \|\bm{y}\|$，因此：
\begin{equation}
\frac{(Q\bm{x})^T (Q\bm{y})}{\|Q\bm{x}\| \|Q\bm{y}\|} = \frac{\bm{x}^T \bm{y}}{\|\bm{x}\| \|\bm{y}\|}
\end{equation}

所以夹角保持不变。

\subsection{性质4.5：反对称矩阵的旋转变换（关键恒等式）}

\textbf{性质}：如果$Q \in SO(3)$（3维特殊正交矩阵），$\bm{v} \in \mathbb{R}^3$是任意向量，$[\bm{v}]$是$\bm{v}$的反对称矩阵，则
\begin{equation}
Q[\bm{v}]Q^T = [Q\bm{v}]
\label{eq:skew_rotation_identity}
\end{equation}

\textbf{反对称矩阵的定义}：对于向量$\bm{v} = (v_x, v_y, v_z)^T \in \mathbb{R}^3$，其反对称矩阵定义为：
\begin{equation}
[\bm{v}] = \begin{bmatrix}
0 & -v_z & v_y \\
v_z & 0 & -v_x \\
-v_y & v_x & 0
\end{bmatrix}
\label{eq:skew_definition}
\end{equation}

\textbf{反对称矩阵的基本性质}：
\begin{itemize}
\item \textbf{反对称性}：$[\bm{v}]^T = -[\bm{v}]$
\item \textbf{叉乘表示}：$[\bm{v}] \bm{u} = \bm{v} \times \bm{u}$（对任意$\bm{u} \in \mathbb{R}^3$）
\end{itemize}

\textbf{证明}：

对于任意$\bm{u} \in \mathbb{R}^3$，考虑$Q[\bm{v}]Q^T$作用在$\bm{u}$上的结果：
\begin{align}
Q[\bm{v}]Q^T \bm{u} &= Q[\bm{v}](Q^T \bm{u}) \\
&= Q(\bm{v} \times (Q^T \bm{u}))
\end{align}

\textbf{关键步骤}：利用叉乘的旋转性质（这可以通过行列式或直接计算证明）：
\begin{equation}
Q(\bm{v} \times \bm{u}) = (Q\bm{v}) \times (Q\bm{u})
\label{eq:cross_product_rotation}
\end{equation}

\textbf{注}：这个性质本身需要证明，但可以通过以下方式验证：
\begin{itemize}
\item 使用行列式性质：$\det([Q\bm{v}], Q\bm{u}, Q\bm{w}) = \det(Q) \det([\bm{v}], \bm{u}, \bm{w})$
\item 或者直接计算验证
\end{itemize}

继续证明：
\begin{align}
Q[\bm{v}]Q^T \bm{u} &= Q(\bm{v} \times (Q^T \bm{u})) \\
&= (Q\bm{v}) \times (Q Q^T \bm{u}) \quad \text{（应用式(\ref{eq:cross_product_rotation})）} \\
&= (Q\bm{v}) \times \bm{u} \quad \text{（因为$Q Q^T = I$）} \\
&= [Q\bm{v}] \bm{u}
\end{align}

由于这对任意$\bm{u} \in \mathbb{R}^3$都成立，我们得到：
\begin{equation}
Q[\bm{v}]Q^T = [Q\bm{v}]
\end{equation}

\textbf{等价表述}：
\begin{equation}
[Q\bm{v}] = Q[\bm{v}]Q^T
\end{equation}

\textbf{几何意义}：
\begin{itemize}
\item 这个恒等式表明：对反对称矩阵进行相似变换（$Q \cdot Q^T$），等价于先对对应的向量进行旋转，再构造反对称矩阵
\item 在机器人学中，这个恒等式在坐标变换中非常重要，特别是在处理角速度、力矩等涉及叉乘的物理量时
\item 这是Lie代数$\mathfrak{so}(3)$到$\mathfrak{so}(3)$的伴随表示（adjoint representation）的基础
\end{itemize}

\textbf{应用}：
\begin{itemize}
\item 用于证明旋转矩阵保持叉乘关系：$(Q\bm{x}) \times (Q\bm{y}) = Q(\bm{x} \times \bm{y})$
\item 在机器人动力学中，用于在不同坐标系间转换角速度和力矩
\item 在SE(3)的伴随表示中起关键作用
\end{itemize}

\subsection{性质5：保持叉乘（3维情况）}

\textbf{性质}：如果$Q \in SO(3)$（3维特殊正交矩阵），$\bm{x}, \bm{y} \in \mathbb{R}^3$，则
\begin{equation}
(Q\bm{x}) \times (Q\bm{y}) = Q(\bm{x} \times \bm{y})
\label{eq:preserve_cross_product}
\end{equation}

\textbf{证明}：我们使用反对称矩阵的性质和性质4.5中的关键恒等式来证明这个重要结果。

\textbf{步骤1}：将叉乘转换为反对称矩阵形式

利用反对称矩阵的性质：$[\bm{v}] \bm{u} = \bm{v} \times \bm{u}$，我们有：
\begin{align}
\bm{x} \times \bm{y} &= [\bm{x}] \bm{y} \\
(Q\bm{x}) \times (Q\bm{y}) &= [Q\bm{x}] (Q\bm{y})
\end{align}

\textbf{步骤2}：应用性质4.5中的关键恒等式

根据性质4.5（式(\ref{eq:skew_rotation_identity})），对于$\bm{x} \in \mathbb{R}^3$，有：
\begin{equation}
[Q\bm{x}] = Q[\bm{x}]Q^T
\end{equation}

\textbf{步骤3}：计算$(Q\bm{x}) \times (Q\bm{y})$
\begin{align}
(Q\bm{x}) \times (Q\bm{y}) &= [Q\bm{x}] (Q\bm{y}) \\
&= (Q[\bm{x}]Q^T)(Q\bm{y}) \quad \text{（应用性质4.5）} \\
&= Q[\bm{x}](Q^T Q)\bm{y} \\
&= Q[\bm{x}] I \bm{y} \quad \text{（因为$Q^T Q = I$）} \\
&= Q[\bm{x}] \bm{y} \\
&= Q(\bm{x} \times \bm{y})
\end{align}

\textbf{因此}：$(Q\bm{x}) \times (Q\bm{y}) = Q(\bm{x} \times \bm{y})$，证毕。

\textbf{几何意义}：
\begin{itemize}
\item 旋转矩阵保持叉乘关系
\item 先旋转两个向量再叉乘，等于先叉乘再旋转结果
\item 在机器人学中，这个性质用于在不同坐标系间转换角速度、力矩等涉及叉乘的物理量
\end{itemize}

\textbf{注}：关键恒等式$Q[\bm{v}]Q^T = [Q\bm{v}]$的详细证明见性质4.5。

\subsection{性质6：列向量和行向量的正交性}

\textbf{性质}：如果$Q = [\bm{q}_1, \bm{q}_2, \ldots, \bm{q}_n]$是正交矩阵，则：
\begin{itemize}
\item 列向量两两正交：$\bm{q}_i^T \bm{q}_j = 0$（$i \neq j$）
\item 列向量是单位向量：$\|\bm{q}_i\| = 1$（$i = 1, 2, \ldots, n$）
\item 行向量也满足相同的性质
\end{itemize}

\textbf{列向量的性质（从$Q^T Q = I$得出）}：

设$Q = [\bm{q}_1, \bm{q}_2, \ldots, \bm{q}_n]$，其中$\bm{q}_i$是第$i$列向量。由$Q^T Q = I$，我们有：
\begin{equation}
Q^T Q = \begin{bmatrix} \bm{q}_1^T \\ \bm{q}_2^T \\ \vdots \\ \bm{q}_n^T \end{bmatrix} [\bm{q}_1, \bm{q}_2, \ldots, \bm{q}_n] = \begin{bmatrix} \bm{q}_i^T \bm{q}_j \end{bmatrix} = I
\end{equation}

因此：
\begin{equation}
\bm{q}_i^T \bm{q}_j = \delta_{ij} = \begin{cases} 1 & \text{如果} i = j \\ 0 & \text{如果} i \neq j \end{cases}
\end{equation}

这意味着：
\begin{itemize}
\item \textbf{正交性}：当$i \neq j$时，$\bm{q}_i^T \bm{q}_j = 0$，即列向量两两正交
\item \textbf{单位长度}：当$i = j$时，$\bm{q}_i^T \bm{q}_i = \|\bm{q}_i\|^2 = 1$，即列向量是单位向量
\end{itemize}

\textbf{行向量的性质（从$Q Q^T = I$得出）}：

设$Q$的行向量为$\bm{r}_1^T, \bm{r}_2^T, \ldots, \bm{r}_n^T$，即：
\begin{equation}
Q = \begin{bmatrix} \bm{r}_1^T \\ \bm{r}_2^T \\ \vdots \\ \bm{r}_n^T \end{bmatrix}
\end{equation}

由$Q Q^T = I$，我们有：
\begin{equation}
Q Q^T = \begin{bmatrix} \bm{r}_1^T \\ \bm{r}_2^T \\ \vdots \\ \bm{r}_n^T \end{bmatrix} \begin{bmatrix} \bm{r}_1 & \bm{r}_2 & \cdots & \bm{r}_n \end{bmatrix} = \begin{bmatrix} \bm{r}_i^T \bm{r}_j \end{bmatrix} = I
\end{equation}

因此：
\begin{equation}
\bm{r}_i^T \bm{r}_j = \delta_{ij}
\end{equation}

这意味着行向量也满足：
\begin{itemize}
\item \textbf{正交性}：当$i \neq j$时，$\bm{r}_i^T \bm{r}_j = 0$，即行向量两两正交
\item \textbf{单位长度}：当$i = j$时，$\bm{r}_i^T \bm{r}_i = \|\bm{r}_i\|^2 = 1$，即行向量是单位向量
\end{itemize}

\textbf{行向量和列向量之间的关系}：

\textbf{重要结论}：对于正交矩阵$Q$，行向量和列向量具有\textbf{完全对称}的性质：
\begin{itemize}
\item 列向量构成标准正交基 $\Longleftrightarrow$ 行向量构成标准正交基
\item 这是因为$Q^T Q = I$和$Q Q^T = I$是等价的（对于方阵）
\item 实际上，$Q$的行向量就是$Q^T$的列向量，而$Q^T$也是正交矩阵
\end{itemize}

\textbf{具体关系}：

设$Q = [\bm{q}_1, \bm{q}_2, \ldots, \bm{q}_n]$，其中$\bm{q}_i$是列向量，则：
\begin{equation}
Q^T = \begin{bmatrix} \bm{q}_1^T \\ \bm{q}_2^T \\ \vdots \\ \bm{q}_n^T \end{bmatrix}
\end{equation}

因此：
\begin{itemize}
\item $Q$的第$i$行向量 = $Q^T$的第$i$列向量 = $\bm{q}_i^T$
\item 由于$Q^T$也是正交矩阵，其列向量（即$Q$的行向量）也构成标准正交基
\end{itemize}

\textbf{矩阵表示}：

如果$Q = [\bm{q}_1, \bm{q}_2, \ldots, \bm{q}_n]$，则：
\begin{align}
Q^T Q &= \begin{bmatrix} \bm{q}_1^T \\ \bm{q}_2^T \\ \vdots \\ \bm{q}_n^T \end{bmatrix} [\bm{q}_1, \bm{q}_2, \ldots, \bm{q}_n] = \begin{bmatrix} \bm{q}_i^T \bm{q}_j \end{bmatrix} = I \\
Q Q^T &= [\bm{q}_1, \bm{q}_2, \ldots, \bm{q}_n] \begin{bmatrix} \bm{q}_1^T \\ \bm{q}_2^T \\ \vdots \\ \bm{q}_n^T \end{bmatrix} = \sum_{i=1}^n \bm{q}_i \bm{q}_i^T = I
\end{align}

第二个等式表明：列向量的外积之和等于单位矩阵，这是标准正交基的另一个特征。

\textbf{总结}：

对于正交矩阵$Q$：
\begin{enumerate}
\item \textbf{列向量}：$\{\bm{q}_1, \bm{q}_2, \ldots, \bm{q}_n\}$构成$\mathbb{R}^n$的标准正交基
\item \textbf{行向量}：$\{\bm{q}_1^T, \bm{q}_2^T, \ldots, \bm{q}_n^T\}$也构成$\mathbb{R}^n$的标准正交基
\item \textbf{对称性}：行向量和列向量的性质完全对称，因为$Q^T$也是正交矩阵
\item \textbf{等价性}：$Q^T Q = I$（列向量标准正交）$\Longleftrightarrow$ $Q Q^T = I$（行向量标准正交）
\end{enumerate}

\subsection{性质7：正交矩阵的乘积}

\textbf{性质}：如果$Q_1$和$Q_2$都是正交矩阵，则$Q_1 Q_2$也是正交矩阵。

\textbf{证明}：
\begin{align}
(Q_1 Q_2)^T (Q_1 Q_2) &= Q_2^T Q_1^T Q_1 Q_2 \\
&= Q_2^T I Q_2 \\
&= Q_2^T Q_2 = I
\end{align}

\textbf{推论}：正交矩阵的逆也是正交矩阵（因为$Q^{-1} = Q^T$）。

\subsection{性质8：特征值的模长}

\textbf{性质}：如果$Q$是正交矩阵，$\lambda$是$Q$的特征值，则
\begin{equation}
|\lambda| = 1
\end{equation}

\textbf{证明}：设$Q\bm{v} = \lambda \bm{v}$，其中$\bm{v} \neq \bm{0}$。则：
\begin{align}
\|Q\bm{v}\|^2 &= \|\lambda \bm{v}\|^2 = |\lambda|^2 \|\bm{v}\|^2 \\
\|Q\bm{v}\|^2 &= \|\bm{v}\|^2 \quad \text{（由性质3）}
\end{align}

因此$|\lambda|^2 = 1$，即$|\lambda| = 1$。

\textbf{几何意义}：正交矩阵的特征值都在单位圆上，表示旋转或反射。

\section{单位正交矩阵的特殊性质}

\subsection{定义回顾}

\textbf{单位正交矩阵}（标准正交矩阵）：正交矩阵$Q$的列向量（或行向量）是标准正交的，即：
\begin{equation}
\bm{q}_i^T \bm{q}_j = \delta_{ij} = \begin{cases} 1 & \text{如果} i = j \\ 0 & \text{如果} i \neq j \end{cases}
\end{equation}

\textbf{注意}：实际上，所有正交矩阵都自动满足这个条件，所以"单位正交矩阵"和"正交矩阵"是等价的。但在某些上下文中，"单位正交矩阵"强调列向量的单位长度性质。

\subsection{单位正交矩阵的额外性质}

\subsubsection{性质1：列向量构成标准正交基}

\textbf{性质}：如果$Q = [\bm{q}_1, \bm{q}_2, \ldots, \bm{q}_n]$是单位正交矩阵，则$\{\bm{q}_1, \bm{q}_2, \ldots, \bm{q}_n\}$构成$\mathbb{R}^n$的标准正交基。

\textbf{证明}：
\begin{itemize}
\item 正交性：$\bm{q}_i^T \bm{q}_j = 0$（$i \neq j$）
\item 单位长度：$\|\bm{q}_i\| = \sqrt{\bm{q}_i^T \bm{q}_i} = 1$
\item 线性无关：由于正交且非零，列向量线性无关
\item 完备性：$n$个线性无关的$n$维向量构成基
\end{itemize}

\subsubsection{性质2：坐标变换的保持性}

\textbf{性质}：如果$\{\bm{q}_1, \bm{q}_2, \ldots, \bm{q}_n\}$是标准正交基，向量$\bm{x}$在这个基下的坐标为$\bm{c} = [c_1, c_2, \ldots, c_n]^T$，则：
\begin{equation}
\bm{x} = \sum_{i=1}^n c_i \bm{q}_i = Q\bm{c}
\end{equation}

且坐标可以通过内积计算：
\begin{equation}
c_i = \bm{q}_i^T \bm{x}
\end{equation}

\textbf{证明}：由于$\bm{q}_i^T \bm{q}_j = \delta_{ij}$，我们有：
\begin{align}
\bm{q}_i^T \bm{x} &= \bm{q}_i^T \sum_{j=1}^n c_j \bm{q}_j \\
&= \sum_{j=1}^n c_j \bm{q}_i^T \bm{q}_j \\
&= \sum_{j=1}^n c_j \delta_{ij} = c_i
\end{align}

\subsubsection{性质3：投影矩阵}

\textbf{性质}：如果$Q$是单位正交矩阵，$Q_r$是$Q$的前$r$列（$r < n$），则投影到$Q_r$的列空间的投影矩阵为：
\begin{equation}
P = Q_r Q_r^T
\end{equation}

\textbf{性质}：
\begin{itemize}
\item $P^2 = P$（幂等性）
\item $P^T = P$（对称性）
\item $\text{rank}(P) = r$
\end{itemize}

\section{具体例子}

\subsection{例子1：2维旋转矩阵}

\textbf{旋转矩阵}：
\begin{equation}
Q = \begin{bmatrix} \cos\theta & -\sin\theta \\ \sin\theta & \cos\theta \end{bmatrix}
\end{equation}

\textbf{验证正交性}：
\begin{align}
Q^T Q &= \begin{bmatrix} \cos\theta & \sin\theta \\ -\sin\theta & \cos\theta \end{bmatrix} \begin{bmatrix} \cos\theta & -\sin\theta \\ \sin\theta & \cos\theta \end{bmatrix} \\
&= \begin{bmatrix} \cos^2\theta + \sin^2\theta & -\cos\theta\sin\theta + \sin\theta\cos\theta \\ -\sin\theta\cos\theta + \cos\theta\sin\theta & \sin^2\theta + \cos^2\theta \end{bmatrix} \\
&= \begin{bmatrix} 1 & 0 \\ 0 & 1 \end{bmatrix} = I
\end{align}

\textbf{验证保持长度}：设$\bm{y} = \begin{bmatrix} y_1 \\ y_2 \end{bmatrix}$，则：
\begin{align}
Q\bm{y} &= \begin{bmatrix} \cos\theta & -\sin\theta \\ \sin\theta & \cos\theta \end{bmatrix} \begin{bmatrix} y_1 \\ y_2 \end{bmatrix} \\
&= \begin{bmatrix} y_1\cos\theta - y_2\sin\theta \\ y_1\sin\theta + y_2\cos\theta \end{bmatrix}
\end{align}

\begin{align}
\|Q\bm{y}\|^2 &= (y_1\cos\theta - y_2\sin\theta)^2 + (y_1\sin\theta + y_2\cos\theta)^2 \\
&= y_1^2(\cos^2\theta + \sin^2\theta) + y_2^2(\sin^2\theta + \cos^2\theta) \\
&\quad + 2y_1 y_2(-\cos\theta\sin\theta + \sin\theta\cos\theta) \\
&= y_1^2 + y_2^2 = \|\bm{y}\|^2
\end{align}

因此$\|Q\bm{y}\| = \|\bm{y}\|$。

\subsection{例子2：3维单位矩阵}

\textbf{单位矩阵}：
\begin{equation}
I = \begin{bmatrix} 1 & 0 & 0 \\ 0 & 1 & 0 \\ 0 & 0 & 1 \end{bmatrix}
\end{equation}

\textbf{性质}：
\begin{itemize}
\item $I^T I = I$，所以$I$是正交矩阵
\item $\det(I) = 1$，所以$I \in SO(3)$
\item 对于任意向量$\bm{y}$，$I\bm{y} = \bm{y}$，所以$\|I\bm{y}\| = \|\bm{y}\|$
\end{itemize}

\subsection{例子3：反射矩阵}

\textbf{反射矩阵}（关于$x$轴的反射）：
\begin{equation}
Q = \begin{bmatrix} 1 & 0 \\ 0 & -1 \end{bmatrix}
\end{equation}

\textbf{验证}：
\begin{itemize}
\item $Q^T Q = I$，所以$Q$是正交矩阵
\item $\det(Q) = -1$，所以$Q$包含反射
\item 对于$\bm{y} = \begin{bmatrix} y_1 \\ y_2 \end{bmatrix}$，$Q\bm{y} = \begin{bmatrix} y_1 \\ -y_2 \end{bmatrix}$
\item $\|Q\bm{y}\|^2 = y_1^2 + y_2^2 = \|\bm{y}\|^2$，长度保持不变
\end{itemize}

\subsection{例子4：数值验证}

\textbf{问题}：验证矩阵$Q = \frac{1}{\sqrt{2}} \begin{bmatrix} 1 & 1 \\ 1 & -1 \end{bmatrix}$是正交矩阵，并验证$\|Q\bm{y}\| = \|\bm{y}\|$，其中$\bm{y} = \begin{bmatrix} 3 \\ 4 \end{bmatrix}$。

\textbf{步骤1}：验证正交性
\begin{align}
Q^T Q &= \frac{1}{2} \begin{bmatrix} 1 & 1 \\ 1 & -1 \end{bmatrix} \begin{bmatrix} 1 & 1 \\ 1 & -1 \end{bmatrix} \\
&= \frac{1}{2} \begin{bmatrix} 2 & 0 \\ 0 & 2 \end{bmatrix} = \begin{bmatrix} 1 & 0 \\ 0 & 1 \end{bmatrix} = I
\end{align}

\textbf{步骤2}：计算$Q\bm{y}$
\begin{align}
Q\bm{y} &= \frac{1}{\sqrt{2}} \begin{bmatrix} 1 & 1 \\ 1 & -1 \end{bmatrix} \begin{bmatrix} 3 \\ 4 \end{bmatrix} \\
&= \frac{1}{\sqrt{2}} \begin{bmatrix} 7 \\ -1 \end{bmatrix}
\end{align}

\textbf{步骤3}：验证长度
\begin{align}
\|\bm{y}\| &= \sqrt{3^2 + 4^2} = \sqrt{9 + 16} = 5 \\
\|Q\bm{y}\| &= \frac{1}{\sqrt{2}} \sqrt{7^2 + (-1)^2} = \frac{1}{\sqrt{2}} \sqrt{49 + 1} = \frac{\sqrt{50}}{\sqrt{2}} = \sqrt{25} = 5
\end{align}

因此$\|Q\bm{y}\| = \|\bm{y}\| = 5$，验证成功！

\section{在机器人学中的应用}

\subsection{应用1：旋转矩阵}

\textbf{旋转矩阵}：在机器人学中，旋转矩阵$R \in SO(3)$是正交矩阵，满足：
\begin{itemize}
\item $R^T R = I$
\item $\det(R) = 1$
\item $\|R\bm{p}\| = \|\bm{p}\|$（保持向量长度）
\end{itemize}

\textbf{意义}：旋转不改变向量的长度，只改变方向。

\subsection{应用2：坐标变换}

\textbf{坐标变换}：在不同坐标系之间变换时，使用正交矩阵：
\begin{equation}
\bm{p}_B = R_{BA} \bm{p}_A
\end{equation}

其中$R_{BA}$是从坐标系$\{A\}$到$\{B\}$的旋转矩阵。

\textbf{性质}：由于$\|R_{BA}\bm{p}_A\| = \|\bm{p}_A\|$，坐标变换保持向量的长度。

\subsection{应用3：雅可比矩阵的QR分解}

\textbf{QR分解}：$J = QR$，其中$Q$是正交矩阵。

\textbf{应用}：在最小二乘问题中，由于$\|Q\bm{y}\| = \|\bm{y}\|$，可以简化计算：
\begin{equation}
\|J\bm{x} - \bm{b}\|^2 = \|QR\bm{x} - \bm{b}\|^2 = \|R\bm{x} - Q^T\bm{b}\|^2
\end{equation}

\subsection{应用4：数值稳定性}

\textbf{优势}：使用正交矩阵进行变换不会放大数值误差，因为：
\begin{itemize}
\item 条件数$\kappa(Q) = 1$（最优条件数）
\item 不会放大舍入误差
\item 数值计算稳定
\end{itemize}

\section{总结}

\subsection{正交矩阵的核心性质}

\begin{enumerate}
\item \textbf{定义}：$Q^T Q = I$（或$Q Q^T = I$）
\item \textbf{逆矩阵}：$Q^{-1} = Q^T$
\item \textbf{行列式}：$\det(Q) = \pm 1$
\item \textbf{保持长度}：$\|Q\bm{y}\| = \|\bm{y}\|$（最重要的性质）
\item \textbf{保持内积}：$(Q\bm{x})^T (Q\bm{y}) = \bm{x}^T \bm{y}$
\item \textbf{保持夹角}：向量之间的夹角不变
\item \textbf{特征值}：$|\lambda| = 1$
\item \textbf{乘积}：正交矩阵的乘积仍然是正交矩阵
\end{enumerate}

\subsection{单位正交矩阵的特殊性}

\begin{itemize}
\item 列向量构成标准正交基
\item 坐标可以通过内积计算：$c_i = \bm{q}_i^T \bm{x}$
\item 投影矩阵形式简单：$P = Q_r Q_r^T$
\end{itemize}

\subsection{在机器人学中的重要性}

\begin{itemize}
\item 旋转矩阵是正交矩阵
\item 坐标变换保持向量长度
\item 数值计算稳定
\item 在QR分解、SVD等算法中广泛应用
\end{itemize}

\section{参考文献}

\begin{itemize}
\item Modern Robotics: Mechanics, Planning, and Control
\item 线性代数教材
\item 矩阵分析相关教材
\end{itemize}

\end{document}